\section{Actual tensor coherence of the induced residue frames}

\begin{lemma}[The induced residue frame of the actual tensor]
\label{lem:native-tc-nef-fs-induced-residue-tensor-frame}
\leanok
\uses{thm:actual-induced-residue-power-frame, lem:cartier-actual-pullback-inverse-tensorator-composition, lem:fc-actual-factor-frame-tensor-naturality}
\lean{EnvelopeProperty.NativeTestConfigurationNefFSTensorResidueFrames.inducedResidueFrame_tensor}
Let $g:Y\to Z$ be the original scheme map and $y$ an original valued
point of $Y$.  Let $\alpha:g^*M\simeq N$ and $\gamma:g^*H\simeq K$
be whole module isomorphisms.  Give the original residue fibres of $M,H$
at $g(y)$ frames $e,d$, and frame the residue fibre of $M\otimes H$ by
the actual inverse pullback tensorator followed by the tensor frame
of $e,d$.  Then its frame induced through the actual tensor comparison
$g^*(M\otimes H)\simeq N\otimes K$ is the inverse residue pullback
tensorator followed by the tensor frame of the two individually induced
original residue frames.
\end{lemma}
\begin{proof}
The induced frame is the inverse pulled original isomorphism, the
original residue pullback-composition comparison and the actual
frame-via-factor construction through the residue field map.
Frame-via-factor tensor coherence and monoidality of actual pullback
composition identify the last two factors on the tensor with their
individual tensor product.  The source inverse tensorator cancels
the mapped source tensorator.  Naturality of the actual inverse
residue pullback tensorator moves the two original isomorphisms to
the individual tensor factors.  Naturality of the whole tensor
frame then gives exactly the claimed induced frame.  Every map,
residue field and isomorphism is the original one.
\end{proof}

\begin{lemma}[The original reference support frames commute with tensor]
\label{lem:native-tc-nef-fs-reference-support-tensor-frame}
\leanok
\uses{lem:fc-actual-factor-frame-tensor-naturality, lem:fc-actual-power-frame-iso-naturality}
\lean{EnvelopeProperty.NativeTestConfigurationNefFSTensorResidueFrames.supportGenericFrame_localPullbackTensor}
For any original map $f:Y\to X$, original open $U\subset X$ with frames
of $L,K$ and original support point $y\in f^{-1}U$, the reference support
frame of the actual tensor frame on $L\otimes K$ equals the tensor of
the two reference support frames.  The common preceding isomorphism is
the pulled inverse tensorator of $f^*$ followed by the inverse residue
pullback tensorator.
\end{lemma}
\begin{proof}
Apply actual frame-via-factor tensor coherence first to the original
local pullback frame and then to its original support-generic factor.
Monoidality of the pullback-composition comparison supplies the
intermediate inverse tensorator, and naturality of the whole frame
construction moves the preceding module comparison through the support
factor.  The two original reference frames are the remaining factors.
\end{proof}
